\documentclass[12pt, a4paper]{article}

% ── Paquetes ────────────────────────────────────────────────────────────────
\usepackage[utf8]{inputenc}
\usepackage[T1]{fontenc}
\usepackage[spanish]{babel}
\usepackage{geometry}
\usepackage{amsmath, amssymb}
\usepackage{graphicx}
\usepackage{booktabs}
\usepackage{array}
\usepackage{multirow}
\usepackage{hyperref}
\usepackage{xcolor}
\usepackage{listings}
\usepackage{caption}
\usepackage{subcaption}
\usepackage{setspace}
\usepackage{parskip}
\usepackage{titlesec}
\usepackage{fancyhdr}
\usepackage{float}

\geometry{
  top=2.5cm, bottom=2.5cm,
  left=3cm,  right=2.5cm
}

\onehalfspacing

% ── Estilo de código ────────────────────────────────────────────────────────
\definecolor{codegray}{rgb}{0.95,0.95,0.95}
\definecolor{codeblue}{rgb}{0.1,0.1,0.6}
\lstset{
  backgroundcolor=\color{codegray},
  basicstyle=\ttfamily\small,
  keywordstyle=\color{codeblue}\bfseries,
  breaklines=true,
  frame=single,
  framesep=4pt,
  numbers=left,
  numberstyle=\tiny\color{gray},
  xleftmargin=1.5em,
}

% ── Encabezado ──────────────────────────────────────────────────────────────
\pagestyle{fancy}
\fancyhf{}
\rhead{\small Inteligencia Artificial 2025--2026}
\lhead{\small \textit{rl-competitive-training-scheduler}}
\cfoot{\thepage}

\hypersetup{
  colorlinks=true,
  linkcolor=codeblue,
  urlcolor=codeblue,
}

% ────────────────────────────────────────────────────────────────────────────
\begin{document}

% ── Portada ─────────────────────────────────────────────────────────────────
\begin{titlepage}
  \centering
  \vspace*{2cm}
  {\Large \textbf{Inteligencia Artificial} \par}
  \vspace{0.3cm}
  {\large Proyecto Final --- Curso 2025--2026 \par}
  \vspace{2cm}
  \rule{\textwidth}{1pt}\\[0.5em]
  {\LARGE \textbf{rl-competitive-training-scheduler} \par}
  \vspace{0.4cm}
  {\large Construcción de rutas de aprendizaje mediante\\
  Deep Q-Network con evaluación pedagógica por LLM \par}
  \rule{\textwidth}{1pt}
  \vspace{3cm}
  {\large Glenda N. Rios Rodríguez \par}
  \vspace{0.2cm}
  {\normalsize Grupo: C312 \par}
  \vspace{0.5cm}
  \vfill
  {\large \today \par}
\end{titlepage}

% ── Tabla de contenidos ─────────────────────────────────────────────────────
\tableofcontents
\newpage

% ════════════════════════════════════════════════════════════════════════════
\section{Descripción del Problema}
% ════════════════════════════════════════════════════════════════════════════
 
En la programación competitiva, la calidad del entrenamiento no depende solo de
la cantidad de problemas resueltos, sino de \textit{qué} problemas se intentan y
en \textit{qué orden}. Un estudiante con 120 minutos disponibles puede malgastar
su sesión en problemas triviales, repetir temáticas ya dominadas, o enfrentar
problemas demasiado difíciles que generan frustración sin aprendizaje real.
 
El problema que aborda este proyecto es el siguiente:
 
\begin{quote}
\textit{Dado un conjunto de $N$ problemas de programación competitiva y un
presupuesto de tiempo $T$ minutos, seleccionar y ordenar una subsecuencia de
problemas que maximice el aprendizaje del estudiante, respetando las restricciones
temporales y adaptándose al estado cognitivo actualizado después de cada intento.}
\end{quote}
 
Este problema presenta tres características que lo hacen especialmente difícil
de resolver con métodos clásicos:
 
\begin{itemize}
  \item \textbf{Secuencialidad}: cada decisión modifica el estado del estudiante
    (ELO por tema, fatiga, tiempo restante), lo que hace que las decisiones futuras
    dependan de las actuales.
  \item \textbf{Incertidumbre}: el resultado de cada intento (resuelto/fallido,
    tiempo empleado) es estocástico y depende del estado interno del estudiante.
  \item \textbf{Objetivos múltiples}: una buena sesión debe equilibrar tasa de
    éxito, cobertura temática, progresión de dificultad y eficiencia temporal.
\end{itemize}
 
El proyecto enmarca este problema como un \textbf{Proceso de Decisión de Markov
(MDP)} y lo resuelve con un agente de aprendizaje por refuerzo basado en
Deep Q-Network (DQN), al que se añade un módulo de evaluación pedagógica mediante
un modelo de lenguaje grande (LLM) local.
 
% ════════════════════════════════════════════════════════════════════════════
\section{Modelado Formal}
% ════════════════════════════════════════════════════════════════════════════
 
\subsection{Proceso de Decisión de Markov}
 
El sistema se modela como un MDP finito por episodio con la tupla
$\mathcal{M} = \langle \mathcal{S},\, \mathcal{A},\, \mathcal{T},\, \mathcal{R},\,
\gamma \rangle$, donde:
 
\paragraph{Espacio de estados $\mathcal{S}$.}
El estado $s_t \in \mathbb{R}^{24}$ en el paso $t$ codifica el estado completo del
estudiante:
\[
  s_t = \bigl[\,r_{\text{norm}},\; f_t,\; t_{\text{used,norm}},\; \text{sr}_t,\;
        \text{ELO}_1^{\text{norm}},\, \ldots,\, \text{ELO}_{20}^{\text{norm}}\,\bigr]
\]
donde $r_{\text{norm}}$ es el rating global normalizado, $f_t \in [0,1]$ la
fatiga acumulada, $t_{\text{used,norm}}$ el tiempo consumido relativo al
presupuesto, $\text{sr}_t$ la tasa de éxito hasta el momento, y
$\text{ELO}_k^{\text{norm}}$ el ELO normalizado para cada uno de los 20 temas
canónicos.
 
\paragraph{Espacio de acciones $\mathcal{A}$.}
La acción $a_t \in \{0, 1, \ldots, N-1\}$ es el índice del problema a intentar
en el paso $t$, sobre un pool fijo de $N = 500$ problemas. Las acciones
inválidas (problemas ya intentados o cuyo tiempo estimado supera el tiempo
restante) se enmascaran con $-\infty$ en el cómputo de Q-valores, garantizando
que el agente solo seleccione acciones factibles.
 
\paragraph{Función de transición $\mathcal{T}$.}
La transición $\mathcal{T}(s_{t+1} \mid s_t, a_t)$ está definida implícitamente
por el modelo del estudiante. La probabilidad de resolución del problema $p$
dado el estado del estudiante sigue un modelo ELO logístico:
\[
  P_{\text{solve}}(p \mid s_t) = \sigma\!\left(\frac{r_{\text{eff}}(s_t) - r_p}{\theta}\right)
\]
donde $\sigma$ es la función sigmoide, $r_p$ es el rating del problema,
$r_{\text{eff}}(s_t) = \min_k \,\text{ELO}_k$ sobre los temas del problema
(hipótesis del eslabón más débil), y $\theta = 400$ es el parámetro de escala
ELO estándar.
 
Si el intento es exitoso, los ELO por tema se actualizan:
\[
  \Delta\,\text{ELO}_k = C_{\text{elo}} \cdot (1 - P_k) \cdot \text{challenge}_k
  \cdot (1 - f_t),
  \quad |\Delta\,\text{ELO}_k| \leq 30
\]
La fatiga se incrementa proporcionalmente al tiempo empleado, con
$\lambda_{\text{fatigue}} = 1.2$.
 
\paragraph{Función de recompensa $\mathcal{R}$.}
La recompensa $r_t = \mathcal{R}(s_t, a_t, s_{t+1})$ tiene dos componentes:
 
\textbf{Recompensa inmediata} por intento:
\[
  r_{\text{step}} =
  \begin{cases}
    r_{\text{éxito}} + \overline{\Delta\text{ELO}} + r_{\text{topic}} +
    r_{\text{challenge}} + r_{\text{prog}} + r_{\text{stag}} +
    r_{\text{overuse}} + r_{\text{div}} & \text{si resuelto} \\
    r_{\text{fracaso}} & \text{si fallido}
  \end{cases}
\]
 
donde los componentes se definen como sigue. Los valores por defecto se
indican entre paréntesis.
 
\textit{Componentes fijos:}
\begin{itemize}
  \item $r_{\text{éxito}} = +10$; \quad $r_{\text{fracaso}} = -2$.
  \item $\overline{\Delta\text{ELO}}$: media de los incrementos de ELO sobre
    los temas del problema (ver ec.\ de actualización de ELO).
  \item $r_{\text{topic}} = +2$ si el problema cubre al menos un tema
    canónico nuevo (no visto en la sesión); $0$ en caso contrario.
\end{itemize}
 
\textit{Bonificación por desafío} $r_{\text{challenge}}$
(peso $w_c = 4.0$), en función del gap $g = r_p - r_{\text{global}}$:
\[
  r_{\text{challenge}} =
  \begin{cases}
    -4.0 & \text{si } g < -300 \quad \text{(trivial)} \\
    -1.5 & \text{si } g > 600  \quad \text{(inalcanzable)} \\
    w_c \cdot \dfrac{g}{200}
      & \text{si } 0 \le g \le 200 \\
    w_c \cdot \max\!\left(0,\; 1 - \dfrac{g - 200}{400}\right)
      & \text{si } 200 < g \le 600 \\
    w_c \cdot 0.2 \cdot \left(1 + \dfrac{g}{300}\right)
      & \text{si } -300 \le g < 0
  \end{cases}
\]
 
\textit{Bonificación por progresión} $r_{\text{prog}}$
(peso $w_p = 5.0$), en función del incremento de dificultad respecto al
problema anterior $\delta = r_p - r_{p-1}$:
\[
  r_{\text{prog}} =
  \begin{cases}
    0 & \text{si es el primer problema} \\
    w_p & \text{si } \delta \ge 300 \\
    w_p \cdot \dfrac{\delta}{300} & \text{si } 0 \le \delta < 300 \\
    \max(-2,\; \delta / 300) & \text{si } \delta < 0
  \end{cases}
\]
 
\textit{Penalización por estancamiento} $r_{\text{stag}} = -6.0$ si el
problema actual y los últimos $W-1$ problemas ($W=3$) son todos triviales
($r_{\text{global}} - r_p > 300$); $0$ en caso contrario.
 
\textit{Penalización por sobreuso de tema} $r_{\text{overuse}}$: para cada
tema $k$ del problema cuya frecuencia relativa $f_k = n_k / n_{\text{tot}}$
supere el umbral $\tau = 0.30$:
\[
  r_{\text{overuse}} = \sum_{k:\, f_k > \tau}
    (-10) \cdot \frac{f_k - \tau}{1 - \tau}
\]
 
\textit{Bonificación por diversidad} $r_{\text{div}}$: para cada tema $k$
del problema con $f_k < \tau$ (poco practicado):
\[
  r_{\text{div}} = \frac{1}{|\text{tags}|}\sum_{k:\, f_k < \tau}
    1.5 \cdot \frac{\tau - f_k}{\tau}
\]
Ambas penalizaciones/bonificaciones se activan solo a partir del tercer intento.
 
\textbf{Recompensa terminal} al finalizar el episodio:
\[
  r_{\text{terminal}} = 10 \cdot \rho_s(\text{ratings}) + 2 \cdot |\mathcal{C}|
  + 0.5 \cdot t_{\text{remaining}}
\]
donde $\rho_s$ es la correlación de Spearman entre el orden de intento y el
rating de los problemas ($+1$ = perfecto fácil$\to$difícil, $-1$ = orden
inverso), $|\mathcal{C}|$ es el número de temas algorítmicos canónicos
distintos cubiertos, y $t_{\text{remaining}}$ es el tiempo restante de sesión
en minutos.
 
\paragraph{Factor de descuento $\gamma$.}
Se usa $\gamma = 0.99$, que valora los refuerzos futuros casi tanto como
los inmediatos, incentivando estrategias a largo plazo dentro del episodio.
 
\subsection{Representación de la acción para la red Q}
 
La red Q no recibe directamente el índice de acción, sino el vector concatenado:
\[
  \mathbf{x}_{t,i} = \bigl[\,s_t \;\|\; \mathbf{p}_i(s_t)\,\bigr] \in \mathbb{R}^{48}
\]
donde $\mathbf{p}_i(s_t) \in \mathbb{R}^{24}$ codifica el problema $i$ en
relación con el estado actual del estudiante:
\[
  \mathbf{p}_i(s_t) = \bigl[\,r_{i,\text{norm}},\; \text{gap}_{\text{norm}},\;
  P_{\text{solve}}(p_i \mid s_t),\; t_{i,\text{norm}},\;
  \text{onehot}(\text{tags}_i)\,\bigr]
\]
El agente calcula $Q(s_t, a_i) = f_\theta(\mathbf{x}_{t,i})$ para todos los
problemas disponibles y selecciona $a_t = \arg\max_i Q(s_t, a_i)$ sobre
acciones válidas.


% ════════════════════════════════════════════════════════════════════════════
\section{Descripción del Dataset}
% ════════════════════════════════════════════════════════════════════════════

\subsection{Fuentes de datos}

Se utilizan dos datasets públicos de Codeforces disponibles en Kaggle:

\begin{table}[H]
\centering
\caption{Datasets utilizados}
\begin{tabular}{llrp{6cm}}
\toprule
\textbf{Dataset} & \textbf{Fuente} & \textbf{Tamaño} & \textbf{Contenido} \\
\midrule
\texttt{codeforces\_dataset.csv} & Kaggle (immortal3) & 14 MB &
  Contest ID, nombre del problema, enunciado, tags con rating embebido (\texttt{*XXXX}) \\
\texttt{codeforces\_problemset.csv} & Kaggle (erchhh) & 20 MB &
  Enunciado, ejemplos de entrada/salida, límite de tiempo, tags en formato lista Python \\
\bottomrule
\end{tabular}
\end{table}

Ambos datasets se combinan mediante un proceso de unión por \texttt{problem\_id}
(construido como concatenación de contest ID y letra del problema, e.g.\ ``325A'').

\subsection{Pipeline de preprocesamiento}

El pipeline transforma los datos crudos en un conjunto de 500 problemas
utilizables por el entorno de simulación:

\begin{enumerate}
  \item \textbf{Construcción de \texttt{problem\_id}}: concatenación de
    \texttt{contest\_id} + letra del problema.
  \item \textbf{Extracción de rating}: parseo del token \texttt{*XXXX} embebido
    en la cadena de tags. Este token es la única fuente de dificultad
    numérica en ambos datasets.
  \item \textbf{Limpieza de tags}: eliminación de \texttt{*XXXX} de la lista
    de tags; normalización a lista de strings en minúsculas.
  \item \textbf{Filtrado}: descarte de filas sin rating o sin enunciado.
  \item \textbf{Muestreo estratificado}: selección de 500 problemas
    distribuidos por bandas de dificultad, con distribución sesgada hacia
    ratings accesibles para reproducir la distribución real del uso del
    sistema.
\end{enumerate}

\begin{table}[H]
\centering
\caption{Distribución por banda de dificultad en el dataset procesado}
\begin{tabular}{lrr}
\toprule
\textbf{Rango de rating} & \textbf{Proporción} & \textbf{Problemas (N=500)} \\
\midrule
800 -- 1199 (Newbie)          & 25\% & 125 \\
1200 -- 1599 (Pupil)          & 30\% & 150 \\
1600 -- 1999 (Specialist)     & 25\% & 125 \\
2000 -- 2399 (Expert)         & 12\% &  60 \\
2400 -- 2799 (Candidate Master) & 5\% & 25 \\
2800+ (Master / Grandmaster)   & 3\% & 15 \\
\bottomrule
\end{tabular}
\end{table}

\subsection{Perfiles de estudiante}

El entorno requiere estudiantes simulados con distribuciones de habilidades
realistas. Se generaron 500 perfiles mediante el LLM local (descrito en la
Sección~\ref{sec:llm}), con ratings globales entre 800 y 2000 y distribuciones
de ELO por tema correlacionadas según el arquetipo generado. Los perfiles se
almacenan en \texttt{data/processed/student\_profiles.json} para garantizar
reproducibilidad de los experimentos.

Los perfiles cubren arquetipos como: especialista en matemáticas débil en grafos,
especialista en DP, estudiante principiante, programador competitivo avanzado, etc.

% ════════════════════════════════════════════════════════════════════════════
\section{Diseño del Algoritmo}
% ════════════════════════════════════════════════════════════════════════════

\subsection{Deep Q-Network (DQN)}

Se implementa el algoritmo DQN estándar~\cite{mnih2015} con las siguientes
características:

\paragraph{Arquitectura de la red Q.}
Un perceptrón multicapa (MLP) con la siguiente topología:
\[
  \mathbb{R}^{48} \xrightarrow{\text{Linear}} \mathbb{R}^{128} \xrightarrow{\text{ReLU}}
  \mathbb{R}^{64} \xrightarrow{\text{ReLU}} \mathbb{R}^{1}
\]
La salida es el Q-valor escalar para el par $(s_t, a_i)$. Se mantienen dos
instancias con la misma arquitectura: la \textit{red online} (actualizada en
cada paso) y la \textit{red target} (sincronizada periódicamente), lo que
estabiliza el entrenamiento al desacoplar los targets de los parámetros
entrenados.

\paragraph{Política $\varepsilon$-greedy.}
\[
  a_t = \begin{cases}
    \text{acción aleatoria válida} & \text{con probabilidad } \varepsilon \\
    \arg\max_{i \text{ válido}} Q_\theta(s_t, a_i) & \text{en caso contrario}
  \end{cases}
\]
$\varepsilon$ decae linealmente de 1.0 a 0.05 durante los primeros 80\% de los
episodios de entrenamiento.

\paragraph{Experience replay.}
Se almacenan transiciones $(s_t, a_t, r_t, s_{t+1}, \text{mask}_{t+1}, \text{done})$
en un buffer circular FIFO de capacidad 10,000. Cada paso de actualización
muestrea aleatoriamente un batch de tamaño 64.

Por eficiencia de memoria, la matriz de problemas
$\mathbf{P}(s_{t+1}) \in \mathbb{R}^{500 \times 24}$ no se almacena en el buffer
(supondría $\sim$240 MB a plena capacidad), sino que se reconstruye en el momento
de entrenamiento desde \texttt{ObservationBuilder}.

\paragraph{Actualización.}
El objetivo de la ecuación de Bellman con red target es:
\[
  y_t = r_t + \gamma \cdot \max_{a' \text{ válido}} Q_{\theta^-}(s_{t+1}, a')
  \cdot (1 - \text{done}_t)
\]
La pérdida es MSE: $\mathcal{L} = \mathbb{E}\left[(y_t - Q_\theta(s_t, a_t))^2\right]$,
minimizada con Adam (lr $= 10^{-3}$). Se aplica gradient clipping con
$\|\nabla\| \leq 1.0$ para mayor estabilidad.

\subsection{Algoritmos de referencia (baselines)}

Se implementaron dos baselines deterministas que comparten la interfaz
abstracta \texttt{ProblemSelector}:

\paragraph{Greedy.}
En cada paso selecciona el problema disponible que maximiza la ratio:
\[
  \text{score}(p) = \frac{P_{\text{solve}}(p) \cdot \Delta\text{ELO}_{\text{esperado}}(p)}{t_{\text{estimado}}(p)}
\]
Es un heurístico adaptativo: reacciona a los cambios de ELO y fatiga tras cada
intento.

\paragraph{Knapsack 0/1.}
Al inicio del episodio resuelve el problema de la mochila 0/1 con programación
dinámica en $O(N \cdot W)$, donde $W$ es el presupuesto de tiempo discretizado
en minutos. El valor de cada problema es:
\[
  V(p) = r_p + 50 \cdot |\{k \in \text{tags}(p) : k \notin \text{topics\_seen}\}|
\]
Una vez elegido el conjunto óptimo, los problemas se ejecutan en orden ascendente
de dificultad. Es el único baseline \textit{no adaptativo}: toma todas las
decisiones al inicio y no reacciona a resultados intermedios.

(El algoritmo Rollout, explorado en experimentos preliminares, mostró un
rendimiento muy inferior con recompensas medias negativas y un coste
computacional inasumible, por lo que se descartó de la comparativa final.)

% ════════════════════════════════════════════════════════════════════════════
\section{Rol del LLM en el Sistema}
\label{sec:llm}
% ════════════════════════════════════════════════════════════════════════════

El sistema integra un LLM local (\texttt{llama3.2:3b} vía Ollama) en dos
roles funcionalmente distintos y complementarios.

\subsection{Rol 1: Generación de perfiles de estudiante realistas}

\textbf{Problema que resuelve.}
Un perfil de estudiante generado aleatoriamente (con ELOs uniformes e
independientes por tema) no es representativo de estudiantes reales, en
los que las habilidades están fuertemente correlacionadas: un estudiante con ELO
alto en DP tiende a tener ELO razonable en matemáticas, mientras que quien es
fuerte en grafos generalmente domina también DFS y árboles.

\textbf{Mecanismo.}
Se consulta al LLM con un prompt que solicita un perfil de estudiante de
programación competitiva con un arquetipo específico. El LLM responde en
formato JSON estructurado con ELOs por tema que reflejan las correlaciones
aprendidas de su preentrenamiento sobre datos de la comunidad competitiva:

\begin{lstlisting}[language=Python, caption=Estructura del perfil generado por el LLM]
{
  "archetype": "Math specialist struggling with graphs",
  "global_rating": 1350,
  "session_budget_min": 120,
  "topic_ratings": {
    "math": 1800, "number theory": 1700,
    "dp": 1200, "graphs": 900, ...
  }
}
\end{lstlisting}

Si Ollama no está disponible, el sistema cae en arquetipos predefinidos
(\textit{fallback} determinista), garantizando que los experimentos sean
reproducibles independientemente del LLM.

\textbf{Impacto en el sistema.}
Perfiles realistas producen condiciones de evaluación más representativas:
el agente DQN y los baselines son evaluados sobre una distribución de
estudiantes más cercana a la real, lo que da mayor validez externa a la
comparación experimental.

\subsection{Rol 2: Evaluación pedagógica de sesiones completadas}

\textbf{Problema que resuelve.}
Las métricas cuantitativas (recompensa total, Spearman, cobertura de temas)
no capturan la calidad pedagógica de una secuencia desde una perspectiva
didáctica. Dos sesiones con la misma recompensa pueden diferir
significativamente en su valor educativo.

\textbf{Mecanismo.}
Tras completar una sesión, se construye un prompt que describe al LLM la
secuencia completa de problemas intentados (ID, rating, tags, resultado,
recompensa) y el estado final del estudiante. El LLM actúa como ``coach''
experto y devuelve una evaluación estructurada en JSON:

\begin{lstlisting}[caption=Estructura de la evaluación pedagógica del LLM]
{
  "score": 8,
  "difficulty_progression": "good",
  "topic_diversity": "average",
  "challenge_level": "appropriate",
  "explanation": "La sesion mostro una progresion ...",
  "recommendation": "Incluir mas problemas de grafos ..."
}
\end{lstlisting}

La puntuación (1--10) y las dimensiones cualitativas (\textit{difficulty
progression}, \textit{topic diversity}, \textit{challenge level}) complementan
las métricas numéricas con una interpretación en lenguaje natural, accesible
para el usuario final.

\textbf{Por qué el LLM es necesario aquí.}
La evaluación pedagógica requiere razonamiento sobre conceptos abstractos
(``¿esta secuencia es apropiada para este estudiante?'') que no pueden
reducirse a una función escalar simple sin perder información semántica
relevante. El LLM proporciona esta capacidad de evaluación cualitativa de
forma directa, sin necesidad de diseñar y validar una rúbrica manual.

% ════════════════════════════════════════════════════════════════════════════
\section{Metodología Experimental}
% ════════════════════════════════════════════════════════════════════════════

\subsection{Entrenamiento del agente DQN}

El agente se entrenó durante 8{,}000 episodios sobre el pool de 500 problemas,
con perfiles de estudiante rotados secuencialmente del conjunto de 500
perfiles preprocesados (cada perfil se reutiliza 16 veces a lo largo del
entrenamiento). Este número de episodios permite que $\varepsilon$ alcance
su valor mínimo en las primeras 6{,}400 iteraciones, dejando las últimas
1{,}600 en régimen casi puramente explotador. Parámetros de entrenamiento:

\begin{table}[H]
\centering
\caption{Hiperparámetros de entrenamiento}
\begin{tabular}{ll}
\toprule
\textbf{Parámetro} & \textbf{Valor} \\
\midrule
Episodios de entrenamiento    & 8{,}000 \\
$\varepsilon_{\text{start}}$  & 1.0 \\
$\varepsilon_{\text{end}}$    & 0.05 \\
Decaimiento de $\varepsilon$  & lineal, 80\% de episodios (6{,}400 ep.) \\
Learning rate (Adam)          & $10^{-3}$ \\
Batch size                    & 64 \\
Capacidad del buffer          & 10{,}000 \\
Sincronización de target      & cada 10 episodios \\
$\gamma$                      & 0.99 \\
Gradient clipping             & 1.0 \\
\bottomrule
\end{tabular}
\end{table}

\subsection{Protocolo de evaluación}

Se realizaron experimentos con $n = 50$ episodios por algoritmo sobre la misma
colección de 50 perfiles de estudiante, garantizando la comparabilidad de los
resultados. Los algoritmos evaluados son DQN, Knapsack y Greedy. Las métricas
registradas en cada episodio son:

\begin{itemize}
  \item \textbf{Recompensa total}: suma de recompensas inmediatas del episodio.
  \item \textbf{Progresión de Spearman} ($\rho_s$): correlación de Spearman
    entre el orden de intento y el rating del problema ($+1$ = perfecto
    fácil$\to$difícil).
  \item \textbf{Cobertura temática}: número de temas algorítmicos canónicos
    distintos cubiertos.
  \item \textbf{Tasa de éxito}: problemas resueltos / intentados.
  \item \textbf{Tiempo usado minutos}: tiempo total consumido en la sesión.
  \item \textbf{Mejora de rating}: ELO global final $-$ ELO global inicial.
\end{itemize}

La significancia estadística se evalúa con el test de Mann-Whitney U (no
paramétrico, apropiado dada la distribución no normal de las métricas)
con nivel de significancia $\alpha = 0.05$.

% ════════════════════════════════════════════════════════════════════════════
\section{Resultados y Análisis}
% ════════════════════════════════════════════════════════════════════════════

Los experimentos se realizaron con 50 perfiles de estudiante diferentes (50
episodios) para cada uno de los tres selectores principales: DQN, Knapsack y
Greedy. El algoritmo Rollout se descartó en esta fase final por su bajo
rendimiento en experimentos preliminares. Las Tablas~\ref{tab:resumen}
y~\ref{tab:comparativa} resumen las métricas obtenidas.

\begin{table}[H]
\centering
\caption{Resultados agregados: DQN, Knapsack y Greedy (50 episodios cada uno)}
\label{tab:resumen}
\resizebox{\textwidth}{!}{%
\begin{tabular}{lcccccc}
\toprule
\textbf{Selector} & \textbf{Recompensa} & \textbf{Spearman} & \textbf{Cobertura} &
\textbf{Tasa éxito} & \textbf{Tiempo usado} & \textbf{Mejora rating} \\
 & \textit{media (sd)} & \textit{media (sd)} & \textit{media (sd)} &
 \textit{media (sd)} & \textit{media (sd)} & \textit{media (sd)} \\
\midrule
DQN      & \textbf{38.17} (21.99) & 0.422 (0.404) & \textbf{4.16} (1.98) & \textbf{0.568} (0.213) & 119.92 (28.05) & 48.99 (133.58) \\
Knapsack & 31.89 (21.04) & \textbf{0.602} (0.324) & 3.10 (1.58) & 0.460 (0.188) & 119.39 (28.34) & \textbf{56.40} (134.80) \\
Greedy   & 16.60 (19.94) & 0.065 (0.438) & 6.30 (1.74) & 0.435 (0.177) & 118.95 (28.27) & 35.59 (122.30) \\
\bottomrule
\end{tabular}%
}
\end{table}

\subsection{Análisis de la recompensa total}

El DQN obtiene la recompensa media más alta (38.17), superando al Knapsack
(31.89) en 6.28 puntos y al Greedy (16.60) en más del doble. La diferencia
entre DQN y Greedy es estadísticamente significativa (Mann-Whitney U,
$p < 0.0001$). Aunque la diferencia entre DQN y Knapsack no alcanza el umbral
de significancia convencional ($p = 0.172$), la ventaja numérica del DQN es
consistente a lo largo de los 50 episodios.

\subsection{Progresión de dificultad (Spearman)}

El Knapsack mantiene la mejor progresión de dificultad (0.602), muy superior
a la del DQN (0.422) y a la prácticamente nula del Greedy (0.065). Las
diferencias son significativas: DQN vs. Knapsack ($p = 0.008$) y DQN vs.
Greedy ($p < 0.0001$). El orden ascendente de dificultad del Knapsack es
estructuralmente óptimo para esta métrica, pero no se traduce en una mejor
recompensa total.

\subsection{Cobertura temática}

El Greedy cubre la mayor cantidad de temas distintos (6.30), seguido por el
DQN (4.16) y finalmente el Knapsack (3.10). Sin embargo, la alta cobertura
del Greedy viene acompañada de una progresión de dificultad deficiente y
una recompensa total baja. El DQN equilibra cobertura y progresión, mientras
que el Knapsack sacrifica cobertura para maximizar la progresión. Todas las
diferencias son significativas ($p < 0.01$).

\subsection{Tasa de éxito}

El DQN presenta la mayor tasa de éxito (0.568), superando al Knapsack (0.460)
y al Greedy (0.435). La diferencia entre DQN y Knapsack es significativa
($p = 0.009$) y también entre DQN y Greedy ($p = 0.001$). Esto indica que el
DQN aprende a seleccionar problemas que el estudiante virtual tiene una
probabilidad realista de resolver, sin caer en la sobreexplotación de
problemas triviales (como sucedía con el Rollout en experimentos previos).

\subsection{Mejora de rating}

El Knapsack obtiene la mayor mejora media de rating (56.40), seguido del DQN
(48.99) y el Greedy (35.59). Las diferencias no son estadísticamente
significativas debido a la alta varianza de esta métrica (desviación típica
$\sim$130 puntos). La mejora de rating está fuertemente influenciada por el
perfil inicial del estudiante y por la aleatoriedad del modelo de éxito/fracaso.

\subsection{Análisis cualitativo por tipos de perfil}

El DQN destaca especialmente en perfiles de estudiante asimétricos o con
puntos débiles muy marcados. Por ejemplo, en el arquetipo
``Math specialist struggling with graphs'' (episodio 2), el DQN alcanza una
recompensa de 62.28 frente a 0.32 del Knapsack y 24.65 del Greedy. En
``Advanced Data Structures and Algorithms Specialist'' (episodio 30), el DQN
obtiene 82.87, superando ampliamente al Knapsack (50.17) y al Greedy (37.85).

Por el contrario, en perfiles muy homogéneos o con todos los ELOs equilibrados,
el Knapsack puede ser competitivo o ligeramente superior, pero la ventaja del
DQN en escenarios realistas (habilidades no uniformes) respalda su elección
como algoritmo principal para el sistema.

\subsection{Síntesis de los resultados}

Los experimentos con 50 episodios por algoritmo confirman que el DQN aprende
una política que maximiza la recompensa total, combinando una tasa de éxito
alta con una cobertura temática intermedia y una progresión de dificultad
aceptable. El Knapsack es el baseline más competitivo, especialmente en
progresión de dificultad, pero su naturaleza estática lo limita ante perfiles
de estudiante heterogéneos. El Greedy, aunque cubre muchos temas, carece de
progresión y obtiene la peor recompensa. Estos resultados validan el enfoque
de aprendizaje por refuerzo profundo para el problema de selección óptima de
problemas bajo restricciones de tiempo.

% ════════════════════════════════════════════════════════════════════════════
\section{Limitaciones y Posibles Mejoras}
% ════════════════════════════════════════════════════════════════════════════

\subsection{Limitaciones actuales}

\paragraph{Modelo del estudiante simplificado.}
El \texttt{StudentModel} es una aproximación probabilística, no un modelo
validado empíricamente. Las constantes de fatiga, actualización de ELO y
tiempo de resolución se calibraron manualmente. Un modelo entrenado con datos
reales de Codeforces (tiempos de envío, historial de aceptaciones) produciría
simulaciones más fieles.

\paragraph{Ausencia de Double DQN y Prioritized Replay.}
La implementación usa DQN estándar. Extensiones conocidas como Double DQN
(reduce sobreestimación de Q-valores) y Prioritized Experience Replay
(muestrea transiciones más informativas) podrían mejorar la convergencia
y la calidad de la política aprendida.

\paragraph{Pool de problemas fijo.}
El pool de 500 problemas es estático. En un sistema real, el pool debería
actualizarse periódicamente con nuevos problemas de Codeforces. La
generalización del agente a problemas no vistos durante el entrenamiento
no ha sido evaluada.

\paragraph{Evaluación LLM sin validación humana.}
Las puntuaciones pedagógicas del LLM no han sido comparadas con evaluaciones
de instructores humanos. La validez del LLM como evaluador pedagógico es
una hipótesis razonable pero no verificada en este trabajo.

\paragraph{Varianza alta en rating improvement.}
La desviación estándar de la mejora de rating ($\sim$110--130 puntos) es
muy elevada en todos los algoritmos, lo que indica que esta métrica está
fuertemente influenciada por el perfil inicial del estudiante y la
aleatoriedad del modelo. Se necesitan más episodios o estratificación por
rating inicial para obtener estimaciones más precisas.

\subsection{Posibles mejoras}

\begin{enumerate}
  \item \textbf{Double DQN / Dueling DQN}: incorporar estas variantes para
    reducir la sobreestimación de Q-valores y mejorar la estabilidad del
    entrenamiento.
  \item \textbf{Modelo del estudiante basado en datos reales}: entrenar los
    parámetros del modelo con historiales de resolución reales de Codeforces.
  \item \textbf{LLM como función de recompensa aumentada}: en lugar de
    usar el LLM solo como evaluador \textit{post-hoc}, integrar su puntuación
    pedagógica directamente en la función de recompensa durante el
    entrenamiento (reward shaping).
  \item \textbf{Pool dinámico de problemas}: conectar el sistema a la API
    de Codeforces para actualizar el pool con problemas recientes.
  \item \textbf{Generalización a otros jueces en línea}: extender el preprocesador
    para soportar datasets de LeetCode o AtCoder.
  \item \textbf{Validación con usuarios reales}: comparar las sesiones
    recomendadas por el DQN con sesiones diseñadas por instructores expertos.
\end{enumerate}

% ════════════════════════════════════════════════════════════════════════════
% Referencias
% ════════════════════════════════════════════════════════════════════════════
\begin{thebibliography}{9}

\bibitem{mnih2015}
V. Mnih et al.,
``Human-level control through deep reinforcement learning,''
\textit{Nature}, vol. 518, pp. 529--533, 2015.

\bibitem{sutton2018}
R. S. Sutton and A. G. Barto,
\textit{Reinforcement Learning: An Introduction}, 2nd ed.
MIT Press, 2018.

\bibitem{elo1978}
A. E. Elo,
\textit{The Rating of Chessplayers, Past and Present}.
Arco Publishing, 1978.

\bibitem{gymnasium}
M. Towers et al.,
``Gymnasium,''
\textit{arXiv:2407.17032}, 2024.

\bibitem{ollama}
Ollama,
``Run large language models locally,''
\url{https://ollama.com}, 2024.

\bibitem{codeforces}
Codeforces,
``Competitive programming judge,''
\url{https://codeforces.com}.

\end{thebibliography}

\end{document}